\section{Conclusion}
\label{sec:conclusion}

We presented a systematic analysis of the dimensionality of belief space, generating \nbeliefs diverse belief statements and analyzing their covariance structure through three complementary PCA methods. Our key finding is that belief space is dramatically lower-dimensional than its surface diversity: 9 principal components explain 50\% of attitudinal variance, 93 explain 90\%, and the conditional influence structure requires only 25 dimensions for 90\% coverage. The dominant axis (28.6\% of variance) is a clear progressive--conservative dimension, with subsequent components capturing institutional trust, individualism, spirituality, and gender attitudes.

These results have three practical implications. First, approximately 25 well-chosen diagnostic questions can reconstruct 90\% of someone's belief profile, informing efficient survey design. Second, the low-dimensional value structure encoded by LLMs can be explicitly extracted and audited, contributing to AI alignment efforts. Third, our five-component model provides a richer framework for understanding opinion structure than the traditional two-axis models in political science.

Future work should validate these components against human survey data from the World Values Survey and General Social Survey, test whether dimensionality is consistent across different LLMs, and explore nonlinear methods (autoencoders, UMAP) that may reveal additional structure beyond linear PCA. The question of whether ``bridge beliefs'' connect otherwise independent dimensions---and whether targeting them could efficiently shift opinion clusters---remains an open and consequential direction.
